\section{Introduction}\label{sec:E-introduction}

The critical activity fixes the weight of each honeycomb polygon, but
the sum over disjoint polygons surrounding a point requires control
on every spatial scale. We determine
the growth of this nesting partition function by first calculating a
partition function for polygons separating two prescribed points on a
cylinder. Two independent pressure estimates then relate the cylinder
to planar nests. The same analysis also gives bounds for large polygon
counts and for the length of critical self-avoiding paths in a strip.

\subsection{The model and its exponents}

We realize the honeycomb lattice as the dual of the equilateral
triangular tiling with side length one and a horizontal edge direction.
A triangular band has height $\sqrt3/2$. A polygon is an unoriented
simple cycle of the dual graph. Its length $|P|$ is its number of
vertices, and its critical weight is $\rho^{|P|}$, where
\[
 \rho=\frac1{\sqrt{2+\sqrt2}}.
\]
Nienhuis predicted this activity through the dilute $O(n)$ model
\cite{Nienhuis1982}; Duminil-Copin and Smirnov proved that its
reciprocal is the honeycomb connective constant \cite{DCS2012}.
Their parafermionic observable supplies the local boundary cancellation
used here. We also use the finite strip transfer construction and the
critical strip-crossing estimate of \cite{StripCompanion}.

Identify the plane with $\mathbb C$, with the horizontal lattice
direction real and a face center at zero. For even $N=2m$, quotient by
the translation $iN\sqrt3/2$ and mark the images of $0$ and
$m e^{2\pi i/3}$.
Their vertical separation is $m$ bands. This is the balanced cylinder
geometry used throughout: a lattice cut follows $m$ edges in direction
$e^{2\pi i/3}$ and then $m$ in direction $e^{\pi i/3}$, closing after
one period. Compactifying the cylinder's two infinite ends gives a
sphere. A polygon \emph{separates the markers} if they lie in distinct
components of its complement on that sphere. Such a polygon may wind
around the cylinder, or may be contractible and enclose exactly one
marker in its disk.

Give every separating polygon an additional factor $\chi$, called its
loop fugacity, and define
\[
 Z_N(\chi)=\sum_{\mathcal P}\prod_{P\in\mathcal P}
                                  \chi\rho^{|P|}.
\]
The sum runs over finite families of pairwise vertex-disjoint separating
polygons, including the empty family with weight one. The transfer
construction in Section~\ref{sec:02-cylinder} proves convergence of
these masses. Every separating polygon meets the finite arc of the
prescribed cut between the markers, so the number of disjoint polygons
is bounded at fixed $N$. Thus $Z_N$ is a polynomial with nonnegative
coefficients and $Z_N(0)=1$.

\begin{theorem}[Cylinder fugacity exponent]
\label{thm:04-cylinder-exponent}\label{E:cylinder-exponent}
For every fixed $\chi>0$, along even $N$,
\[
 Z_N(\chi)=N^{\tau(\chi)+o(1)},\qquad
 \tau(\chi)=
 \begin{cases}
 \displaystyle\frac16-\frac{2}{3\pi^2}
              \bigl(\arccos(\chi/2)\bigr)^2,&0<\chi\le2,\\[4pt]
 \displaystyle\frac16+\frac{2}{3\pi^2}
              \bigl(\operatorname{arcosh}(\chi/2)\bigr)^2,&\chi\ge2.
 \end{cases}
\]
\end{theorem}

For the planar counterpart, let $P_\chi(r)$ be the sum of
$\prod_{P\in\mathcal P}\chi\rho^{|P|}$ over disjoint nests around
one fixed face center, with every polygon of diameter at most $r$;
again include the empty nest. Theorem~\ref{thm:05-nest-normalization}
proves the bounded-factor comparison
\[
 P_\chi(c_0N)^2\le Z_N(\chi)
       \le C_{\chi,c_0}P_\chi(c_0N)^2
\]
for every fixed $\chi\ge0$ and sufficiently small fixed $c_0>0$.
Consequently, for every fixed $\chi>0$,
\[
 P_\chi(r)=r^{\tau(\chi)/2+o(1)}.
\]
The square has a geometric origin. Removing essential polygons and
large contractible polygons costs a bounded factor; the remaining
polygons form two independent small planar nests, one near each marker.
Figure~\ref{fig:E-planar-comparison} shows the two planar neighborhoods.
The two pressure calculations below provide the estimates needed for
this removal.

In particular, $Z_N(2)=N^{1/6+o(1)}$ and
$P_2(r)=r^{1/12+o(1)}$. Fugacity two arises in the boundary winding
identity because an enclosing polygon has two orientations. If $S_k$
is the same fugacity-two nesting mass around a middle face of a strip
of $2k$ bands, with no diameter restriction, then
Corollary~\ref{E:strip-nesting} gives
\[
 S_k=k^{1/12+o(1)}.
\]

\subsection{Uniform counts and path length}

The fixed-fugacity theorem concerns one parameter at a time. To control
polygon counts that grow with $N$, we prove a separate estimate uniform
over an unbounded fugacity range. Define the critical coefficient mass
$w_N(a)$ by
\[
 Z_N(\chi)=\sum_{a\ge0}w_N(a)\chi^a.
\]
Thus $w_N(a)$ sums $\prod_P\rho^{|P|}$ over families with exactly
$a$ separating polygons. Section~\ref{sec:06-all-count} proves
\[
 Z_N(2\cosh v)\le\exp\!\bigl(C\log N(1+v^2)\bigr)
       \quad(v\ge0),
\]
with $C$ independent of $v$, and hence
\[
 w_N(a)\le\exp\!\left(C\log N-c\frac{a^2}{\log N}\right)
       \quad(a\ge0).
\]
These are unnormalized coefficient masses. The uniform fugacity bound
uses the coarse real-sector estimate before any counts are truncated;
the fixed-parameter asymptotic alone would not give it.

The nesting law also controls a path-length observable. A \emph{port}
is the midpoint of a triangular-tile edge. In a strip of $h$ bands,
consider self-avoiding paths from a fixed bottom boundary port to
either boundary, with no other boundary visit. They include their
initial and final half-edges, and $n(\gamma)$ counts the visited dual
vertices. Proposition~\ref{E:raw-first-moment} proves
\[
 \sum_\gamma\rho^{n(\gamma)}n(\gamma)
       \le h^{13/12+o(1)}.
\]
Under the normalized critical weights, their expected length conditional
on reaching the top boundary of the bottom $\lfloor h/4\rfloor$ bands is
$h^{4/3+o(1)}$ (Proposition~\ref{prop:02-length-reduction}).
The difference between these exponents comes from the conditioning
event, whose unnormalized mass is comparable to $h^{-1/4}$ by the
strip-crossing theorem of \cite[Theorem~1.1]{StripCompanion}.
A positive identity relating visits to an interior edge to nests
around its two adjacent faces supplies the required length comparison.

\subsection{Proof strategy}

The proof has an algebraic and analytic route to the cylinder exponent,
followed by a geometric route to nesting. We describe the two routes
separately because the pressure estimates do not follow from the
cylinder asymptotic.

\paragraph{A finite representation with its normalization.}
Section~\ref{sec:02-cylinder} constructs the cylinder vacuum from a
polynomial open-strip state. A vacuum state records the pairings of
occupied ports made by paths on one side of a cut, with no external
source. Contracting the two half-cylinder states with the separation
rule gives $Z_N$. Degree bounds and reductions at special parameters
determine this contraction exactly.

Writing $\chi=2\cos\eta$, a finite contour deformation expresses
its numerator as a sum of particle integrals with factor $e^{i\eta M}$.
The integer $M$ fixes the two groups' elementary particle totals as
$m-M$ and $m+M$. Each group contains main particles and
residue-generated far single particles and fused pairs, the latter
counted twice in the elementary degree.
Only terms of elementary degree $N$ enter the identity.

\paragraph{The logarithmic contribution of the particle integrals.}
Section~\ref{sec:03-gas-mean} centers each main group at the mean
density $m\varrho(u)$, where $\varrho$ is a fixed probability density
with exponential tails. Walls at distance proportional to $\log N$
leave only $O(1)$ mean mass outside; this is the spatial scale of the
particle calculation. The real quadratic energy of the remaining fluctuations
controls their imbalance, their far-particle counts, and excursions
beyond that interval. A Legendre orthogonal-polynomial ensemble gives
a positive reference law for the two large interior groups. Subtracting
its logarithmic interaction removes the diagonal singularity, so the
remaining energy is meaningful on empirical particle measures.

The coefficient of $\log N$ must be kept throughout this comparison.
The relative fluctuation determinant contributes $7/8$, the two
reference partition masses contribute $-1/2$, and the scalar vacuum
denominator contributes $-5/24$. Their sum is
\[
 \frac78-\frac12-\frac5{24}=\frac16.
\]
To obtain the quadratic term in $\eta$, the phase $e^{i\eta M}$
must be retained during integration.
Section~\ref{sec:E-reference} proves the reference integration
bounds. Section~\ref{sec:04-gas-fluctuations} realizes the needed
complex displacement by moving actual particles and computing its
Jacobian. Because the interior particle numbers also change with $M$,
Section~\ref{sec:E-interpolation} constructs a holomorphic approximation
to the integer-dimensional integrals before moving the imbalance sum.
The combined particle and imbalance-contour deformations convert the
phase into a quadratic contribution to the real action.
These steps yield, locally uniformly for $|\operatorname{Re}\eta|<\pi$,
\[
 \frac{\log|Z_N(2\cos\eta)|}{\log N}
 \le \frac16-\frac{2}{3\pi^2}\operatorname{Re}(\eta^2)+o(1).
\]
Section~\ref{sec:E-cylinder-exponent} obtains the matching lower
exponent from this harmonic upper envelope. The anchor $\eta=\pi/2$
has the exact value $Z_N(0)=1$; subharmonicity propagates near equality,
and coefficient positivity transfers it to each fixed positive
fugacity. Thus no lower estimate of a single complex sector is needed.

\paragraph{From the cylinder to planar polygons and strip paths.}
Section~\ref{sec:05-pressure} varies the weight of essential loops
in an annular transfer. Its first two pressure derivatives at zero fugacity control
both their total mass and their first horizontal-span moment. These
bounds make the cost of removing separating essential polygons uniform
in $N$.
Section~\ref{sec:05-open-pressure} instead varies the weight of all
loops in an open strip. Its first pressure derivative at zero fugacity
is affine in $N$ up to an error $O(N^{-1})$. Counting each shape's
vertical placements and comparing two strip widths then gives
\[
 \sum_{[P]:\,\diam P>r}\rho^{|P|}\le C(1+r)^{-2}.
\]
Here $[P]$ is an unrooted, unoriented polygon modulo triangular-lattice
translations. This estimate controls large contractible polygons.
Section~\ref{sec:05-nesting} uses both removal bounds to separate
the two planar nests, then proves the length--nest identity and the
raw and conditioned strip-length estimates.

\subsection{Earlier work}

The electric exponent has a direct Coulomb-gas predecessor. Gamsa and
Cardy study the change in weight of loops separating marked points
\cite[Section~3, equation~(6)]{GamsaCardy2006}. In the dilute
zero-loop specialization, their scaling dimension for
$\chi=2\cos\eta$, $0<\chi\le2$, is
$x=\eta^2/(3\pi^2)-1/12$, giving the predicted two-point power
$-2x=\tau(\chi)$. Their argument assumes the Gaussian-field scaling
description. Cardy's annulus calculation likewise permits an independent
weight for winding loops and gives the corresponding radial power
\cite{Cardy2006Annulus}. The present argument establishes the finite
lattice normalization and the estimates required for the cylinder and
planar observables above; it also proves the stated positive-fugacity
branch for $\chi>2$.

Periodic dilute-loop transfer matrices with seams that change wrapping
loop weights were analyzed by Zhou and Batchelor
\cite{ZhouBatchelor1997}. Their spectral calculation provides related
electric-parameter formulas. The observable here is a contraction with
two specified markers, and includes contractible polygons surrounding
one marker. The finite contraction and its nonvanishing scalar
normalization must therefore be identified in this geometry.
Ikhlef and Cardy related local parafermionic equations to the
Yang--Baxter weights of dilute-loop models \cite{IkhlefCardy2009}.
Glazman and Manolescu proved invariance of the boundary two-point
function for columnwise rhombic half-planes with angles in
$[\pi/3,2\pi/3]$ \cite{GM2019}. We use their local-weight convention,
specialized to the regular honeycomb lattice in the final statements.

Polynomial exchange equations and special-rapidity reductions have
earlier uses in integrable loop constructions. Di Francesco and
Zinn-Justin use them to prove a multiparameter sum rule for the dense
$O(1)$ model \cite{DiFrancescoZinnJustin2005}. Garbali and Nienhuis
develop polynomial states and fusion recurrences for the dilute
$O(1)$ open strip, and obtain determinant expressions for its
normalization \cite{GarbaliNienhuis2017Ground,GarbaliNienhuis2017Sum}.
They also note that the polynomial-state method applies at zero loop
weight \cite[Section~4]{GarbaliNienhuis2017Ground}. Here the
open-strip polynomial state is supplied by \cite{StripCompanion};
we construct its periodic limit, the normalized vacuum quotient,
and the two-marker contraction with their physical interpretation.

The analytic reference measure is a classical orthogonal-polynomial
ensemble. Andr\'eief's identity gives its Gram normalization and
tilted factorial determinants \cite{TracyWidom1998}. The circle
comparison is the finite upper bound associated with the strong
Szeg\H{o} theorem \cite{Simon2005Szego}; its Schur-polynomial proof
is related to the unitary trace calculations of Diaconis--Shahshahani
and Diaconis--Evans \cite{DiaconisShahshahani1994,DiaconisEvans2001}.
We include these arguments in the required
normalization. To retain the logarithmic exponent, the present
application also needs uniform source bounds for growing physical
tests, a particle realization of the complex quadratic change of
variables, and continuation of the integer dimensions. These are
the roles of Sections~\ref{sec:E-reference}--\ref{sec:E-interpolation}.
